\newpage
\section {Relatório de Agosto}

\subsection{Propostas para o mês:}
\begin{enumerate}
    \item Criação do menu inicial:
    \begin{itemize}
        \item Ter a opção de iniciar um novo jogo.
        \item Criar um botão para as opções, créditos e para sair do jogo.
    \end{itemize}
    
    \item Criação das artes:
    \begin{itemize}
        \item Estilizar o menu inicial.
        \item Criar as artes das primeiras cartas do jogo, o suficiente para uma partida rápida. 
    \end{itemize}
    
    \item Implementação das mecânicas básicas: (pontos de vida, tabuleiro, baralho e ações do jogador)
    \begin{itemize}
        \item Implementar os pontos de vida do jogador (encerrando a partida quando reduzidos à zero).
        \item Implementar o tabuleiro, com os espaços para posicionar as cartas.
        \item Programar o funcionamento do baralho.
        \item Programar as ações disponíveis para os jogadores.
    \end{itemize}
\end{enumerate}

\subsection{Atividades Realizadas:}

\subsubsection{Engenharia do projeto}

Após todas as discussões dos últimos meses, uma documentação inicial do projeto através de diagramas foi feita para organizar as classes do programa. Também foi feito um Kanban para submeter as tarefas que foram e estão sendo feitas, com base nos estudos de engenharia de software. Todos os diagramas e documentos importantes foram submetidos em um Drive compartilhado com os membros, bem como no Github

\subsubsection{Diretório do Github e Menu Inicial}

Com tudo pronto para iniciar a programação, o diretório do projeto foi criado com uma estrutura inicial para acomodar os principais componentes do código, como as bibliotecas, classes, interfaces, etc. Essa estrutura foi feita para guiar os membros quanto ao modelo que os arquivos devem seguir para manter a coerência e organização. Além disso, o arquivo README também documenta como devem ser os padrões de commit e como o programa deve ser compilado, dando o passo a passo para que seja de fácil entendimento.

As bibliotecas e funções do Raylib também foram integradas, assim criando a base do menu que será implementado nas próximas semanas, mas que no momento pode ser inicializado fazendo o jogo “rodar” em uma tela preta. É um bom primeiro passo, pois demonstra que as principais funções de criação de janela do Raylib estão funcionando. 


\subsubsection{Design das cartas}

Os responsáveis pelo design do jogo também criaram um documento com as principais ideias de cada carta, desde a parte mecânica até a parte visual. Algumas artes de treino foram feitas, seguindo uma temática futurista que será a identidade visual do jogo. Para os próximos passos, prosseguirão para a estilização do menu e criação das cartas que irão entrar para essa versão atual do jogo.

\subsubsection{Metas futuras}

Mesmo que as implementações em si do jogo ainda não estejam completas, considerando o que era previsto para este mês, mantendo a organização, essas metas logo serão alcançadas no próximo mês (Setembro). 